\subsubsection{Tensão do efeito de viga}

A tensão do efeito de viga é dada pela equação:

\begin{equation}
  \sigma_v = \frac{n_e \cdot T \cdot L_t}{\pi \cdot {D_t}^2 \cdot h_1}
  \label{eq:tensao_devido_ao_efeito_de_viga}
\end{equation}

Onde:

\begin{itemize}
  \item $\sigma_v$: tensão do efeito de viga em \si{\mega\pascal}.
  \item $n_e$: número de entradas (2 entradas).
  \item $T$: força de tração no cabo em \si{\newton} \eqref{eq:forca_de_tracao_no_cabo}.
  \item $L_t$: comprimento do tambor em \si{\meter} \eqref{eq:comprimento_do_tambor}.
  \item $D_t$: diâmetro do tambor em \si{\meter} (Tabela \ref{tab:dados_do_tambor_adotado}).
  \item $h_1$: espessura remanescente sob a ranhura em \si{\meter} (Figura \ref{fig:dim_tambor}).
\end{itemize}

\vspace{1em}

Substituindo os valores na equação \eqref{eq:tensao_devido_ao_efeito_de_viga}:

\begin{align*}
  \sigma_v & = \frac{2 \cdot 25377{,}8 \cdot 2{,}112}{\pi \cdot 0{,}34794^2 \cdot 0{,}01} \\
  \sigma_v & = \boxed{\SI{28{,}19}{\mega\pascal}}
\end{align*}

\subsubsection{Tensão de flexão local}

A tensão de flexão local é dada pela equação:

\begin{equation}
  \sigma_f = 0{,}96 \cdot T \sqrt[4]{\frac{1}{{D_t}^2 \cdot {h_1}^6}}
  \label{eq:tensao_de_flexao_local}
\end{equation}

Onde:

\begin{itemize}
  \item $\sigma_f$: tensão de flexão local em \si{\mega\pascal}.
  \item $T$: força de tração no cabo em \si{\newton} \eqref{eq:forca_de_tracao_no_cabo}.
  \item $D_t$: diâmetro do tambor em \si{\meter} (Tabela \ref{tab:dados_do_tambor_adotado}).
  \item $h_1$: espessura remanescente sob a ranhura em \si{\meter} (Figura \ref{fig:dim_tambor}).
\end{itemize}

\vspace{1em}

Substituindo os valores na equação \eqref{eq:tensao_de_flexao_local}:

\begin{align*}
  \sigma_f & = 0{,}96 \cdot 25377{,}8 \sqrt[4]{\frac{1}{0{,}34794^2 \cdot 0{,}01^6}} \\
  \sigma_f & = \boxed{\SI{41{,}3}{\mega\pascal}}
\end{align*}

\subsubsection{Tensão de compressão do enrolamento do cabo}

A tensão de compressão do enrolamento do cabo é dada pela equação:

\begin{equation}
  \sigma_{c_e} = \frac{0{,}5 \cdot T \cdot L_t}{p \cdot h_1 + 0{,}112 \cdot p^2}
  \label{eq:tensao_de_compressao_do_enrolamento_do_cabo}
\end{equation}

Onde:

\begin{itemize}
  \item $\sigma_{c_e}$: tensão de compressão do enrolamento do cabo em \si{\mega\pascal}.
  \item $T$: força de tração no cabo em \si{\newton} \eqref{eq:forca_de_tracao_no_cabo}.
  \item $L_t$: comprimento do tambor em \si{\meter} \eqref{eq:comprimento_do_tambor}.
  \item $p$: passo da ranhura em \si{\meter} (Tabela \ref{tab:dimensoes_ranhuras}).
  \item $h_1$: espessura remanescente sob a ranhura em \si{\meter} (Figura \ref{fig:dim_tambor}).
\end{itemize}

\vspace{1em}

Substituindo os valores na equação \eqref{eq:tensao_de_compressao_do_enrolamento_do_cabo}:

\begin{align*}
  \sigma_{c_e} & = \frac{0{,}5 \cdot 25377{,}8 \cdot 2{,}112}{0{,}021 \cdot 0{,}01 + 0{,}112 \cdot 0{,}021^2} \\
  \sigma_{c_e} & = \boxed{\SI{103{,}32}{\mega\pascal}}
\end{align*}

\subsubsection{Tensão resultante}

A tensão resultante é dada pela equação:

\begin{equation}
  \sigma_{\text{res}} = \sqrt{{\left(\sigma_v + \sigma_f\right)}^2 + {\sigma_{c_e}}^2}
  \label{eq:tensao_resultante}
\end{equation}

Onde:

\begin{itemize}
  \item $\sigma_{\text{res}}$: tensão resultante em \si{\mega\pascal}.
  \item $\sigma_v$: tensão do efeito de viga em \si{\mega\pascal}.
  \item $\sigma_f$: tensão de flexão local em \si{\mega\pascal}.
  \item $\sigma_{c_e}$: tensão de compressão do enrolamento do cabo em \si{\mega\pascal}.
\end{itemize}

\vspace{1em}

Substituindo os valores na equação \eqref{eq:tensao_resultante}:

\begin{align*}
  \sigma_{\text{res}} & = \sqrt{{\left(28{,}19 + 41{,}3\right)}^2 + {103{,}32}^2} \\
  \sigma_{\text{res}} & = \boxed{\SI{124{,}52}{\mega\pascal}}
\end{align*}

Verificando a tensão resultante com a tensão admissível (\eqref{eq:tensao_admissivel_tambor}):

\begin{align*}
  \sigma_{\text{res}}       & \le \sigma_{\text{adm}}                             \\
  \SI{124{,}52}{\mega\pascal} & \le \SI{148{,}21}{\mega\pascal} \rightarrow \text{OK}
\end{align*}
